% =====================================================================
%  Capítulo 1. Introducción  (Bloque UMA obligatorio)
%  ESTADO: andamiaje + objetivos/contribuciones estables. La redacción
%  definitiva de §1.1 y §1.6 se cerrará tras consolidar resultados.
% =====================================================================
\chapter{Introducción}
\label{cap:introduccion}

% --- Objetivo del capítulo ------------------------------------------
% Presentar contexto, motivación, objetivos, preguntas de investigación
% y estructura del documento.

\section{Contexto y motivación}
\label{sec:contexto}
\pendienteredaccion{} La complejidad cognitiva es una métrica de
mantenibilidad que estima el esfuerzo humano necesario para comprender un
fragmento de código \pendientecita{}. A diferencia de la complejidad
ciclomática, penaliza especialmente el anidamiento de estructuras de
control, lo que convierte a los condicionales anidados en uno de sus
principales contribuyentes \pendientecita{}. Esta sección situará la
frecuencia de dichos condicionales en código real (apoyándose en el
escaneo masivo descrito en la \autoref{sec:corpus}) y motivará por qué
su simplificación es relevante para la calidad del software.

\section{Problema abordado}
\label{sec:problema-intro}
El trabajo aborda la \emph{combinación segura de sentencias condicionales
anidadas} en Java, con foco en el patrón canónico
\lstinline|if (A) { if (B) { S } }| $\rightarrow$ \lstinline|if (A && B) { S }|.
La descripción detallada del problema se ofrece en el
\autoref{cap:problema}.

\section{Objetivos}
\label{sec:objetivos-intro}
\paragraph{Objetivo general.}
Diseñar, implementar y evaluar un enfoque para detectar y refactorizar
sentencias condicionales anidadas en Java con el fin de reducir la
complejidad cognitiva del código, preservando su semántica cuando
corresponda; y comparar este enfoque con refactorizaciones generadas por
grandes modelos de lenguaje bajo un protocolo experimental controlado.

\paragraph{Objetivos específicos.}
\begin{enumerate}
  \item Identificar y formalizar los casos en que es posible combinar
  condicionales anidados de forma segura, documentando precondiciones de
  aplicabilidad y resultado de la transformación.
  \item Diseñar e implementar un prototipo determinista que detecte dichas
  oportunidades sobre el árbol de sintaxis abstracta (AST) y aplique la
  transformación.
  \item Medir el impacto de las refactorizaciones sobre la complejidad
  cognitiva con una métrica definida y reproducible.
  \item Comparar las refactorizaciones del prototipo con las producidas por
  grandes modelos de lenguaje, bajo protocolo controlado, en términos de
  corrección, completitud y calidad.
\end{enumerate}

\section{Preguntas de investigación}
\label{sec:rq-intro}
% IMPORTANTE: RQ oficiales. Reproducir SIEMPRE literalmente. No reformular.
\begin{description}
  \item[RQ1.] ¿En qué casos se pueden combinar sentencias condicionales anidadas y cuál sería el resultado?
  \item[RQ2.] ¿Qué impacto tienen estas refactorizaciones en la complejidad cognitiva?
  \item[RQ3.] ¿Pueden los grandes modelos de lenguaje realizar esta refactorización de forma correcta y automática?
\end{description}
Su operacionalización (sin reformular el texto oficial) se detalla en la
\autoref{sec:operacionalizacion}.

\section{Contribuciones del TFM}
\label{sec:contribuciones}
\begin{enumerate}
  \item La formalización de un conjunto cerrado de precondiciones
  (P1--P5) para la combinación segura de \lstinline|if| anidados, con su
  catálogo de motivos de descarte trazable.
  \item Un prototipo determinista basado en AST con un pipeline
  experimental reproducible (carga de corpus, detección, transformación,
  medición y exportación de artefactos).
  \item Una comparación controlada del prototipo frente a LLMs bajo un
  protocolo congelado y un oráculo de validación automático.
  \item Un paquete de replicación con corpus, scripts y artefactos
  versionados.
\end{enumerate}

\section{Estructura del documento}
\label{sec:estructura}
\pendienteredaccion{} El \autoref{cap:antecedentes} presenta los
antecedentes; el \autoref{cap:problema} describe el problema y
operacionaliza las RQ; el \autoref{cap:propuesta} detalla la propuesta
(diseño, arquitectura, implementación, diseño experimental, resultados y
amenazas a la validez); y el \autoref{cap:conclusiones} recoge
conclusiones y trabajo futuro. (Texto a afinar al cierre.)
